% Contracapa -- editada à mão (pedido do usuário, 2026-11-10): página inteira no azul institucional da capa (opção D),
% com dados básicos da publicação e a logo Prefeitura Rio | IPP clara no pé, na mesma posição da capa (\capalogo,
% de pretextual/capa.tex). Sem foto, sem número de página. Entra no fim de relatorio.tex, depois dos apêndices.
% A marca d'água "EM DESENVOLVIMENTO" das outras páginas fica por baixo do azul, como na capa.
\clearpage
\thispagestyle{empty}
\begin{tikzpicture}[remember picture, overlay]
  \fill[ippnavy!82!black] (current page.south west) rectangle (current page.north east);
  % fio dourado curto e bloco de dados, alinhados à margem do texto da capa (2,4 cm)
  \fill[c4] ([xshift=2.4cm, yshift=-17.6cm]current page.north west) rectangle ([xshift=3.6cm, yshift=-17.7cm]current page.north west);
  \node[anchor=north west, inner sep=0, align=left, text=white, font={\frauncesdisplay\fontsize{24}{29}\selectfont}]
    at ([xshift=2.4cm, yshift=-18.3cm]current page.north west)
    {Diagnóstico da Primeira Infância Carioca};
  \node[anchor=north west, inner sep=0, align=left, text=white!80!ippnavy, font=\small]
    at ([xshift=2.4cm, yshift=-19.8cm]current page.north west)
    {Relatório técnico \textperiodcentered{} Rio de Janeiro, novembro de 2026\\[0.35cm]
     Instituto Municipal de Urbanismo Pereira Passos (IPP)\\
     Coordenadoria de Pesquisa, Avaliação e Políticas Públicas\\
     Em parceria com a Casa Civil da Prefeitura da Cidade do Rio de Janeiro};
  \capalogo{logo_ipp_clara.png}
\end{tikzpicture}
